\documentclass[11pt,a4paper]{article}

% --- Packages ---
\usepackage[utf8]{inputenc}
\usepackage[T1]{fontenc}
\usepackage{lmodern}
\usepackage[margin=1in]{geometry}
\usepackage{amsmath,amssymb,amsfonts}
\usepackage{graphicx}
\usepackage{booktabs}
\usepackage{multirow}
\usepackage{hyperref}
\usepackage{cleveref}
\usepackage{enumitem}
\usepackage{xcolor}
\usepackage{natbib}
\usepackage{setspace}
\usepackage{array}
\usepackage{colortbl}

\hypersetup{
    colorlinks=true,
    linkcolor=blue!70!black,
    citecolor=green!50!black,
    urlcolor=blue!60!black
}

\setstretch{1.15}

% --- Custom commands ---
\newcommand{\molssd}{\textsc{MolSSD}}
\newcommand{\xt}{\mathbf{x}_t}
\newcommand{\xzero}{\mathbf{x}_0}
\newcommand{\R}{\mathbb{R}}
\newcommand{\E}{\mathbb{E}}
\newcommand{\Loss}{\mathcal{L}}

% --- Title ---
\title{%
  \vspace{-1cm}
  \textbf{Molecular Scale Space Diffusion (MolSSD):} \\[6pt]
  \large A Principled Multi-Scale Framework for Generative Design \\
  of UV-Absorbing Molecular Materials
}

\author{
  Umesh \\
  \textit{Research Proposal} \\[4pt]
  \small Target Venue: \textit{Nature}
}

\date{March 2026}

\begin{document}
\maketitle
\thispagestyle{empty}

% ============================================================
% 1. EXECUTIVE SUMMARY
% ============================================================
\section{Executive Summary}

The design of novel molecular materials with targeted optical properties remains one of the
grand challenges at the intersection of computational chemistry and machine learning.
Current generative diffusion models for molecules operate exclusively at full atomic
resolution, lacking a principled multi-scale theory that would enable both computational
efficiency and chemical interpretability. We propose \textbf{Molecular Scale Space Diffusion
(\molssd{})}, the first framework to rigorously adapt scale-space diffusion
theory~\citep{ssd2026} --- which unifies Gaussian blurring with diffusion-based generation
--- to molecular graphs. \molssd{} defines a \emph{molecular scale space} that
hierarchically coarsens molecules from full atomic detail through functional groups,
fragments, and chromophore skeletons to a centroid representation, with mathematically
grounded non-isotropic posteriors at each resolution-changing step. Applied to the urgent
challenge of discovering novel UV-absorbing materials for protective coatings in autonomous
systems, \molssd{} will be pre-trained on the OMol25 dataset~\citep{omol25} (140M+ DFT
calculations, 83M molecules) and conditioned on absorption spectra, photostability, and
coating compatibility. We expect to generate $>$10,000 candidate UV absorbers, validate the
top candidates through TD-DFT screening and wet-lab synthesis, and deliver both a
foundational theoretical contribution and practical novel materials. This work bridges
scale-space theory from computer vision, coarse-graining from statistical mechanics, and
generative modeling from machine learning into a unified framework for molecular design.


% ============================================================
% 2. INTRODUCTION AND MOTIVATION
% ============================================================
\section{Introduction and Motivation}

\subsection{The Molecular Generation Problem}

Designing molecules with specific, application-driven properties is a central challenge
across pharmaceutical development, materials science, and chemical engineering. The
molecular design space is astronomically large --- estimates place the number of
drug-like molecules alone at $10^{23}$ to $10^{60}$ --- rendering exhaustive enumeration
impossible~\citep{bilodeau2022generative}. Generative models have emerged as a powerful
paradigm for navigating this space, learning the underlying distribution of chemically
valid, synthetically accessible molecules and sampling from it to propose novel candidates
with desired properties~\citep{gomez2018automatic}.

Among generative approaches, diffusion models have recently achieved state-of-the-art
results in molecular generation. Equivariant diffusion models (EDM)~\citep{hoogeboom2022edm}
generate 3D molecular structures respecting the SE(3) symmetry of physical space.
Extensions including GeoDiff~\citep{xu2023geodiff}, GeoLDM~\citep{xu2023geoldm},
MiDi~\citep{vignac2023midi}, GCDM~\citep{morehead2024gcdm}, and
EQGAT-diff~\citep{le2024eqgat} have progressively improved sample quality, training
stability, and property control. Yet all of these models share a fundamental limitation:
they operate at a \emph{single, fixed atomic resolution} throughout the entire diffusion
process.

\subsection{Limitations of Current Diffusion Models}

Working exclusively at full atomic resolution imposes three interrelated costs:

\begin{enumerate}[leftmargin=*]
  \item \textbf{Computational inefficiency.} Every denoising step operates over all $N$
    atoms, even in the early stages of generation when the molecular structure is still
    largely noise. The computational cost scales as $O(N^2)$ or worse per step, making
    generation of large molecules ($N > 100$ heavy atoms) prohibitively slow.

  \item \textbf{Lack of interpretability.} The intermediate states of a standard diffusion
    trajectory are noisy point clouds with no chemically meaningful interpretation. There is
    no ``scaffold first, then decorate'' logic --- a design principle that human chemists
    routinely exploit. Users of generative models cannot inspect or intervene at intermediate
    scales.

  \item \textbf{No multi-scale theory.} Despite decades of scale-space
    theory~\citep{lindeberg1994scalespace,witkin1987scalespace,koenderink1984structure} in
    computer vision and coarse-graining theory~\citep{noid2008cg,kmiecik2016cg} in molecular
    simulation, no unified mathematical framework connects these ideas with diffusion-based
    generation. Existing hierarchical approaches~\citep{qiang2023hierdiff} use ad-hoc
    fragmentation schemes (BRICS~\citep{degen2008brics}, Murcko
    scaffolds~\citep{bemis1996murcko}) that lack information-theoretic justification.
\end{enumerate}

\subsection{The UV-Absorber Materials Challenge}

Ultraviolet radiation causes photodegradation of polymers, coatings, and composite
materials, leading to mechanical failure, discoloration, and loss of functional
properties~\citep{yousif2013photodegradation}. UV-absorbing molecules are incorporated
into protective coatings to mitigate this damage by converting UV photons into harmless
thermal energy~\citep{gugumus1995stabilization,schaller2009uv}.

Current commercial UV absorbers fall into a small number of chemical families ---
benzotriazoles, benzophenones, triazines, and hydroxyphenyltriazines --- each with
well-documented limitations~\citep{bojinov2012photostabilizers,cantrell2012uv}:
\begin{itemize}[leftmargin=*]
  \item \textbf{Narrow absorption windows:} most commercial absorbers cover only portions
    of the UV-A or UV-B range, requiring multi-component blends for broad-spectrum
    protection;
  \item \textbf{Photodegradation:} many absorbers themselves degrade under prolonged UV
    exposure, reducing coating lifetime;
  \item \textbf{Migration and leaching:} small-molecule absorbers can migrate to coating
    surfaces and leach into the environment;
  \item \textbf{Limited tunability:} the structure--property relationships governing
    absorption wavelength, molar absorptivity, and photostability are poorly understood
    outside established scaffolds.
\end{itemize}

These limitations are especially acute for \textbf{autonomous systems} --- self-driving
vehicles, drones, and outdoor robotics --- whose sensor housings, LiDAR windows, and
structural composites face continuous UV exposure in uncontrolled
environments~\citep{akay2024autonomous}. The growing deployment of such systems creates
an urgent need for next-generation UV absorbers with broad-spectrum coverage
(280--400\,nm), exceptional photostability ($>$1000 hours), and compatibility with
modern coating chemistries.

\subsection{Why This Is Timely}

Three recent developments make this project uniquely well-timed:

\begin{enumerate}[leftmargin=*]
  \item \textbf{Scale Space Diffusion theory} (March 2026): The SSD
    paper~\citep{ssd2026} provides the first rigorous mathematical framework unifying
    Gaussian blurring with diffusion models, including the non-isotropic posterior
    derivation essential for resolution-changing steps. This theoretical foundation is
    only weeks old.

  \item \textbf{OMol25 dataset} (2025): The Open Molecules 2025 dataset~\citep{omol25}
    provides 140M+ DFT calculations covering 83M unique molecular systems with
    unprecedented chemical diversity, including the ground-state energetics needed for
    pre-training molecular models at scale.

  \item \textbf{Growing demand for UV protection in autonomous systems:} The rapid
    expansion of autonomous vehicle fleets and drone operations creates immediate
    commercial demand for superior UV-absorbing materials.
\end{enumerate}


% ============================================================
% 3. RESEARCH GAPS
% ============================================================
\section{Research Gaps}

We identify seven specific gaps in the current literature that \molssd{} directly addresses.

\begin{enumerate}[leftmargin=*,label=\textbf{Gap~\arabic*.}]

  \item \textbf{No information-theoretic basis for multi-scale molecular generation.} \\
    Existing hierarchical molecular generation methods~\citep{qiang2023hierdiff} lack a
    formal information-theoretic grounding for how molecular information should be
    distributed across scales. \molssd{} provides this through the rate-distortion
    interpretation of the SSD forward process, where each resolution level corresponds to
    an optimal lossy compression of the molecular representation at a given information
    budget.

  \item \textbf{Ad-hoc molecular fragmentation.} \\
    Current coarsening strategies rely on heuristic fragmentation rules (BRICS
    retrosynthetic cuts~\citep{degen2008brics}, Murcko scaffold
    decomposition~\citep{bemis1996murcko}) that have no optimality guarantees.
    \molssd{} replaces these with spectral graph coarsening derived from the molecular
    graph Laplacian~\citep{chung1997spectral}, providing a principled,
    chemistry-aware decomposition that preserves the most informative structural
    features at each scale.

  \item \textbf{No formal non-isotropic posteriors in molecular coarse-graining.} \\
    When a diffusion model transitions between resolution levels, the reverse-process
    posterior becomes non-isotropic because different degrees of freedom are affected
    differently by the coarsening operator. No existing molecular generative model
    handles this correctly. \molssd{} derives and implements the exact non-isotropic
    posterior using the Lanczos algorithm~\citep{lanczos1950iteration} for efficient
    eigendecomposition of the coarsening operator.

  \item \textbf{Computational inefficiency at full atomic resolution.} \\
    All current 3D molecular diffusion models~\citep{hoogeboom2022edm,morehead2024gcdm}
    expend equal computational effort at every denoising step, regardless of the
    signal-to-noise ratio. \molssd{}'s multi-resolution approach allocates computation
    proportional to the structural detail being resolved: early steps operate on compact
    coarse representations, with full atomic resolution engaged only in the final
    refinement stage.

  \item \textbf{Lack of interpretable intermediate states.} \\
    The intermediate samples $\xt$ in standard molecular diffusion are noisy and
    chemically meaningless. \molssd{}'s scale-space trajectory produces interpretable
    intermediates --- a chromophore skeleton at coarse scale, functional group placement at
    medium scale, full atomic detail at fine scale --- enabling chemist-in-the-loop
    intervention and design insight.

  \item \textbf{No unified theory connecting coarse-graining and generative modeling.} \\
    Physical coarse-graining~\citep{noid2008cg} and generative
    diffusion~\citep{ho2020ddpm,song2021scorebased} have developed as separate fields
    with no formal bridge. \molssd{} provides this bridge by showing that the SSD forward
    process on molecular graphs is equivalent to a systematic coarse-graining, with deep
    connections to the renormalization group~\citep{wilson1983renormalization,kadanoff1966scaling}.

  \item \textbf{No generative design framework for UV-absorbing materials.} \\
    Despite significant commercial interest, no existing generative model has been
    specifically designed or applied to the systematic generation of UV-absorbing
    molecules with targeted spectral properties. \molssd{} fills this gap with
    UV-specific conditioning on absorption wavelength, molar absorptivity,
    photostability, and coating compatibility.

\end{enumerate}


% ============================================================
% 4. PROPOSED METHOD
% ============================================================
\section{Proposed Method}

\subsection{Core Mathematical Framework}

\molssd{} adapts the Scale Space Diffusion (SSD) framework~\citep{ssd2026} from
continuous image domains to discrete molecular graphs. We represent a molecule as a
tuple $\mathcal{M} = (\mathbf{X}, \mathbf{H}, \mathbf{A})$ where
$\mathbf{X} \in \R^{N \times 3}$ are atomic coordinates,
$\mathbf{H} \in \R^{N \times d_h}$ are node features (atom type, charge, hybridization),
and $\mathbf{A} \in \{0,1\}^{N \times N}$ is the adjacency matrix.

\paragraph{Forward process.}
The forward process combines standard Gaussian noise injection with spectral graph
coarsening. At timestep $t$, the noised molecular state is:
\begin{equation}\label{eq:forward}
  \xt = \mathbf{M}_{1:t}\,\xzero + \sigma_t\,\boldsymbol{\epsilon},
  \qquad \boldsymbol{\epsilon} \sim \mathcal{N}(\mathbf{0}, \mathbf{I}),
\end{equation}
where $\mathbf{M}_{1:t} = \mathbf{M}_t \mathbf{M}_{t-1} \cdots \mathbf{M}_1$ is the
cumulative coarsening operator and $\sigma_t$ is the noise schedule. In standard DDPM,
$\mathbf{M}_{1:t} = \sqrt{\bar{\alpha}_t}\,\mathbf{I}$ and the process is purely
noise-driven. In \molssd{}, $\mathbf{M}_t$ is a \emph{spectral coarsening matrix}
derived from the graph Laplacian $\mathbf{L} = \mathbf{D} - \mathbf{A}$ at designated
resolution-changing steps.

\paragraph{Spectral coarsening.}
Given the normalized graph Laplacian $\mathbf{L}_\text{sym} = \mathbf{I} -
\mathbf{D}^{-1/2}\mathbf{A}\mathbf{D}^{-1/2}$ with eigendecomposition
$\mathbf{L}_\text{sym} = \mathbf{U}\boldsymbol{\Lambda}\mathbf{U}^\top$, we define
the coarsening operator at resolution level $r$ as:
\begin{equation}\label{eq:coarsen}
  \mathbf{M}^{(r)} = \mathbf{U}_{:,1:k_r}\,\mathbf{U}_{:,1:k_r}^\top,
\end{equation}
where $k_r < N$ retains only the $k_r$ lowest-frequency eigenvectors, corresponding
to the coarsest structural features of the molecular graph. This is analogous to
Gaussian blurring in the image domain but operates on the graph spectral domain. The
retained eigenvectors naturally capture chemically meaningful groupings: the lowest
frequencies delineate connected components and large-scale topology, intermediate
frequencies resolve functional groups and ring systems, and high frequencies encode
individual atom-level detail.

\paragraph{Molecular scale space.}
We define a hierarchy of five resolution levels that map to chemically interpretable
representations:
\begin{center}
\begin{tabular}{clcl}
  \toprule
  \textbf{Level} & \textbf{Representation} & \textbf{Approx.\ $k_r/N$} & \textbf{Chemical meaning} \\
  \midrule
  $r=0$ & Full atoms          & 1.0   & Complete molecular graph \\
  $r=1$ & Functional groups   & 0.5   & Heavy-atom groups (--OH, --NH$_2$, --COOH, \ldots) \\
  $r=2$ & Fragments           & 0.2   & BRICS-like fragments, ring systems \\
  $r=3$ & Chromophore skeleton& 0.05  & Conjugated backbone, key stereocenters \\
  $r=4$ & Centroid            & $1/N$ & Single representative point \\
  \bottomrule
\end{tabular}
\end{center}

\paragraph{Non-isotropic reverse process.}
At resolution-changing steps in the reverse (generative) direction, the posterior
distribution $q(\mathbf{x}_{t-1} | \xt, \xzero)$ becomes non-isotropic because
the coarsening operator $\mathbf{M}_t$ has non-trivial spectral structure. Following
the SSD derivation, the posterior covariance takes the form:
\begin{equation}\label{eq:posterior}
  \boldsymbol{\Sigma}_{t-1|t} = \left(
    \frac{\mathbf{M}_t^\top \mathbf{M}_t}{\sigma_t^2 - \sigma_{t-1}^2}
    + \frac{\mathbf{I}}{\sigma_{t-1}^2}
  \right)^{-1},
\end{equation}
which is no longer a scalar multiple of the identity. We compute this efficiently
using the Lanczos algorithm~\citep{lanczos1950iteration} to obtain a low-rank
approximation of $\mathbf{M}_t^\top \mathbf{M}_t$, avoiding explicit inversion of
the full $N \times N$ matrix. This is a key theoretical contribution: existing
molecular diffusion models implicitly assume isotropic posteriors, which is only
correct when no resolution change occurs.

\subsection{Molecular Flexi-Net Architecture}

The denoising network $\boldsymbol{\epsilon}_\theta(\xt, t, c)$ must operate at
varying resolutions while maintaining SE(3) equivariance. We propose
\textbf{Molecular Flexi-Net}, a hierarchical E($n$)-equivariant graph neural
network inspired by the Flexi-Net architecture from SSD~\citep{ssd2026} and
adapted to molecular graphs:

\begin{itemize}[leftmargin=*]
  \item \textbf{Resolution-adaptive message passing:} At each resolution level, message
    passing operates on the appropriately coarsened graph. Equivariant graph neural
    network (EGNN) layers~\citep{satorras2021egnn} process coordinates and features
    jointly, preserving SE(3) equivariance through center-of-mass aggregation during
    coarsening.

  \item \textbf{Cross-resolution skip connections:} Feature maps from coarser resolutions
    are projected back to finer grids via the transpose of the coarsening operator,
    providing global context to local refinement steps.

  \item \textbf{Adaptive depth:} Following DiT~\citep{peebles2023dit}, we use adaptive
    layer normalization (adaLN) to condition the network on both the diffusion timestep
    $t$ and the current resolution level $r$, allowing the same parameters to handle
    all scales.

  \item \textbf{SE(3) equivariance throughout:} Coordinate updates use equivariant
    operations exclusively. Coarsening aggregates coordinates via center-of-mass (which
    commutes with rotations and translations), and all message-passing layers use
    radial basis functions on interatomic distances.
\end{itemize}

\subsection{UV-Absorber Conditioning Strategy}

For conditional generation of UV-absorbing molecules, we augment the denoising
network with property conditioning using classifier-free
guidance~\citep{ho2022classifierfree}. The conditioning vector $\mathbf{c}$ encodes:
\begin{itemize}[leftmargin=*]
  \item \textbf{Target absorption spectrum:} A discretized UV-Vis absorption profile
    $A(\lambda)$ over 280--400\,nm, parameterized as $\lambda_\text{max}$, bandwidth,
    and peak molar absorptivity $\varepsilon_\text{max}$;
  \item \textbf{Photostability score:} A scalar encoding predicted resistance to
    photodegradation, derived from TD-DFT excited-state analysis;
  \item \textbf{Coating compatibility:} Solubility parameters (Hansen solubility
    parameters), molecular weight constraints, and functional group requirements for
    covalent incorporation into coating matrices.
\end{itemize}

During training, we randomly drop the conditioning vector with probability $p = 0.1$
to enable classifier-free guidance at inference. At generation time, the guided score
estimate is:
\begin{equation}\label{eq:cfg}
  \tilde{\boldsymbol{\epsilon}}_\theta(\xt, t, \mathbf{c})
  = (1 + w)\,\boldsymbol{\epsilon}_\theta(\xt, t, \mathbf{c})
  - w\,\boldsymbol{\epsilon}_\theta(\xt, t, \varnothing),
\end{equation}
where $w$ is the guidance weight controlling the fidelity--diversity tradeoff.

\subsection{OMol25 Integration}

The Open Molecules 2025 (OMol25) dataset~\citep{omol25} plays three distinct roles in
the \molssd{} pipeline:

\begin{enumerate}[leftmargin=*]
  \item \textbf{Pre-training corpus.} OMol25's 83M unique molecular systems spanning
    diverse chemistries provide an unparalleled pre-training corpus for the Molecular
    Flexi-Net. We pre-train the encoder on a denoising objective over OMol25 geometries
    to learn robust, transferable molecular representations before fine-tuning on
    UV-absorber-specific data.

  \item \textbf{Energy oracle via eSEN.} The eSEN (equivariant Scalable Energy Network)
    models~\citep{fairchem2024} trained on OMol25 provide rapid ground-state energy and
    force predictions for generated candidates, enabling energy-based filtering before
    expensive TD-DFT calculations. This serves as a computationally cheap validity check.

  \item \textbf{Scalability benchmark.} OMol25 contains molecules ranging from small
    organics to systems with $>$100 heavy atoms, making it the ideal testbed for
    demonstrating \molssd{}'s computational scaling advantages over fixed-resolution
    methods on large molecular systems.
\end{enumerate}


% ============================================================
% 5. EXPERIMENTAL PLAN
% ============================================================
\section{Experimental Plan}

\subsection{Phase 1: Foundation and Validation (Months 1--3)}

\paragraph{Objective:} Implement the core \molssd{} framework and validate on standard
benchmarks before tackling UV absorbers.

\begin{itemize}[leftmargin=*]
  \item \textbf{QM9 benchmark}~\citep{ramakrishnan2014qm9}: 134K small molecules
    ($\le$9 heavy atoms). Validate molecular validity, uniqueness, and novelty metrics.
    Compare against EDM~\citep{hoogeboom2022edm}, GeoLDM~\citep{xu2023geoldm}, and
    GCDM~\citep{morehead2024gcdm}. Target: $\ge$95\% validity, competitive FCD scores.

  \item \textbf{GEOM-Drugs benchmark}~\citep{axelrod2022geom}: 304K drug-like
    conformations with molecules up to 181 atoms. Validate scaling behavior and
    conformational accuracy (RMSD to ground truth). Demonstrate wall-clock speedup from
    multi-resolution generation.

  \item \textbf{OMol25 scalability test}: Pre-train on a representative 1M-molecule
    subset of OMol25. Measure training throughput and generation quality as a function
    of molecular size. Quantify the computational savings of multi-resolution generation
    versus full-resolution baselines.

  \item \textbf{Ablation studies}: Systematically evaluate (i) the number of resolution
    levels (2, 3, 5), (ii) spectral versus heuristic coarsening, (iii) isotropic
    versus non-isotropic posteriors, (iv) Min-SNR-5 loss
    weighting~\citep{hang2023minsnr} versus uniform weighting, and (v) ConvexDecay(0.5)
    versus linear resolution schedules.
\end{itemize}

\subsection{Phase 2: UV-Absorber Dataset and Training (Months 3--5)}

\paragraph{Objective:} Curate a UV-absorber dataset and train the conditional \molssd{}
model.

\begin{itemize}[leftmargin=*]
  \item \textbf{Dataset curation:} Assemble a dataset of $\sim$50K known UV-absorbing
    molecules from ChEMBL, PubChem, patent literature, and commercial catalogs.
    For each molecule, compute or extract: (i)~UV-Vis absorption spectrum via
    TD-DFT (CAM-B3LYP/6-31+G(d))~\citep{yanai2004camb3lyp,casida1995tddft},
    (ii)~photostability indicators from excited-state analysis,
    (iii)~synthesizability score~\citep{ertl2009sa},
    (iv)~coating-relevant physicochemical properties.

  \item \textbf{TD-DFT calculations:} Perform TD-DFT calculations on the curated
    dataset using the CAM-B3LYP functional, which is well-validated for
    charge-transfer excited states relevant to UV
    absorption~\citep{yanai2004camb3lyp}. Budget: $\sim$50K single-point
    calculations ($\sim$25K GPU-hours on modern hardware).

  \item \textbf{Conditional model training:} Fine-tune the OMol25-pre-trained
    Molecular Flexi-Net with UV-absorber conditioning. Train with classifier-free
    guidance ($p_\text{drop} = 0.1$). Validate conditional generation quality
    on a held-out test set (10\% of UV-absorber data).
\end{itemize}

\subsection{Phase 3: Generation and Screening (Months 5--7)}

\paragraph{Objective:} Generate a large library of candidate UV absorbers and screen
them computationally.

\begin{itemize}[leftmargin=*]
  \item \textbf{Candidate generation:} Generate 10,000+ candidate molecules conditioned
    on target absorption profiles spanning 280--400\,nm, with guidance weights $w \in
    \{1, 2, 5, 10\}$ to explore the diversity--fidelity tradeoff.

  \item \textbf{eSEN energy filtering:} Use eSEN models~\citep{fairchem2024} pre-trained
    on OMol25 to predict ground-state energies and forces for all candidates. Filter out
    thermodynamically unstable structures (formation energy $> 0.5$\,eV/atom above hull).

  \item \textbf{Synthesizability screening:} Compute SA
    scores~\citep{ertl2009sa} and retrosynthetic
    feasibility~\citep{coley2017retrosynthesis} for energy-stable candidates.
    Retain molecules with SA score $< 4.0$ and at least one viable retrosynthetic route.

  \item \textbf{TD-DFT validation:} Perform full TD-DFT calculations
    (CAM-B3LYP/6-31+G(d,p))~\citep{casida1995tddft} on the top 500 candidates surviving
    energy and synthesizability filters. Compare predicted absorption spectra against
    conditioning targets. Select top 50 candidates for experimental validation.

  \item \textbf{Multi-objective ranking:} Rank final candidates by a composite score
    balancing absorption coverage (280--400\,nm), peak molar absorptivity
    ($\varepsilon_\text{max} > 10^4$\,L\,mol$^{-1}$\,cm$^{-1}$), predicted
    photostability, and synthetic accessibility.
\end{itemize}

\subsection{Phase 4: Wet-Lab Validation (Months 7--9)}

\paragraph{Objective:} Synthesize and experimentally characterize the most promising
computationally-designed UV absorbers.

\begin{itemize}[leftmargin=*]
  \item \textbf{Synthesis:} Select the top 10--20 candidates based on Phase~3 rankings
    and synthetic feasibility assessment. Synthesize via established organic chemistry
    routes, prioritizing molecules with 3--5 step syntheses from commercial starting
    materials.

  \item \textbf{UV-Vis spectroscopy:} Measure experimental absorption spectra in
    solution (multiple solvents) and thin film. Compare against TD-DFT predictions
    and generation targets. Success criterion: $|\lambda_\text{max}^\text{exp} -
    \lambda_\text{max}^\text{pred}| < 20$\,nm for $>$70\% of synthesized compounds.

  \item \textbf{Photostability testing:} Subject synthesized absorbers to accelerated
    UV aging (QUV weatherometer, ASTM G154). Monitor absorption spectrum degradation
    over 500--1000 hours. Target: $<$10\% decrease in peak absorbance after 500 hours.

  \item \textbf{Coating integration tests:} Incorporate top performers into model
    coating formulations (acrylic-urethane and epoxy-amine systems). Evaluate
    UV-protective performance via substrate photodegradation assays. Compare
    against commercial benchmarks (Tinuvin~328, Tinuvin~1130).
\end{itemize}

\subsection{Phase 5: Analysis and Paper (Months 9--11)}

\paragraph{Objective:} Comprehensive analysis, comparison, and manuscript preparation.

\begin{itemize}[leftmargin=*]
  \item \textbf{Theoretical analysis:} Formalize the connection between SSD on molecular
    graphs and renormalization group theory. Prove or empirically verify information
    preservation bounds across resolution levels. Analyze the spectral properties of
    molecular Laplacians and their relationship to chemical coarsening.

  \item \textbf{Benchmarking:} Compile comprehensive comparisons against all baselines
    on QM9, GEOM-Drugs, and the UV-absorber task. Report wall-clock generation times,
    sample quality metrics, and property prediction accuracy.

  \item \textbf{Manuscript preparation:} Write and submit a manuscript to \textit{Nature},
    emphasizing (i)~the theoretical framework, (ii)~the practical discovery of novel UV
    absorbers, and (iii)~the broader implications for multi-scale molecular design.

  \item \textbf{Code and data release:} Open-source the \molssd{} codebase, trained model
    weights, UV-absorber dataset, and generated candidate library.
\end{itemize}


% ============================================================
% 6. TIMELINE
% ============================================================
\section{Timeline}

\begin{table}[h!]
\centering
\small
\renewcommand{\arraystretch}{1.2}
\definecolor{phase1}{HTML}{4472C4}
\definecolor{phase2}{HTML}{ED7D31}
\definecolor{phase3}{HTML}{A5A5A5}
\definecolor{phase4}{HTML}{FFC000}
\definecolor{phase5}{HTML}{5B9BD5}
\newcommand{\cell}[1]{\cellcolor{#1}\hspace{0pt}}
\begin{tabular}{|l|c|c|c|c|c|c|c|c|c|c|c|}
\hline
\textbf{Activity} & \rotatebox{90}{\textbf{M1}} & \rotatebox{90}{\textbf{M2}} & \rotatebox{90}{\textbf{M3}} & \rotatebox{90}{\textbf{M4}} & \rotatebox{90}{\textbf{M5}} & \rotatebox{90}{\textbf{M6}} & \rotatebox{90}{\textbf{M7}} & \rotatebox{90}{\textbf{M8}} & \rotatebox{90}{\textbf{M9}} & \rotatebox{90}{\textbf{M10}} & \rotatebox{90}{\textbf{M11}} \\
\hline
SSD implementation \& QM9         & \cell{phase1} & \cell{phase1} & & & & & & & & & \\
\hline
GEOM-Drugs \& OMol25 scaling      & & \cell{phase1} & \cell{phase1} & & & & & & & & \\
\hline
Ablation studies                   & & & \cell{phase1} & & & & & & & & \\
\hline
UV-absorber dataset curation       & & & \cell{phase2} & \cell{phase2} & & & & & & & \\
\hline
TD-DFT calculations (training)    & & & & \cell{phase2} & \cell{phase2} & & & & & & \\
\hline
Conditional model training         & & & & & \cell{phase2} & & & & & & \\
\hline
Candidate generation (10K)         & & & & & & \cell{phase3} & & & & & \\
\hline
eSEN + SA screening                & & & & & & \cell{phase3} & & & & & \\
\hline
TD-DFT validation (top 500)       & & & & & & \cell{phase3} & \cell{phase3} & & & & \\
\hline
Synthesis (10--20 compounds)       & & & & & & & \cell{phase4} & \cell{phase4} & & & \\
\hline
Spectroscopy \& photostability     & & & & & & & & \cell{phase4} & \cell{phase4} & & \\
\hline
Coating integration tests          & & & & & & & & & \cell{phase4} & & \\
\hline
Theoretical analysis               & & & & & & & & & \cell{phase5} & \cell{phase5} & \\
\hline
Manuscript writing                 & & & & & & & & & & \cell{phase5} & \cell{phase5} \\
\hline
Code/data release                  & & & & & & & & & & & \cell{phase5} \\
\hline
\end{tabular}
\caption{Gantt chart for the 11-month \molssd{} project. Colors indicate phases:
\textcolor{phase1}{$\blacksquare$}~Foundation \& Validation,
\textcolor{phase2}{$\blacksquare$}~UV-Absorber Dataset \& Training,
\textcolor{phase3}{$\blacksquare$}~Generation \& Screening,
\textcolor{phase4}{$\blacksquare$}~Wet-Lab Validation,
\textcolor{phase5}{$\blacksquare$}~Analysis \& Paper.}
\label{tab:timeline}
\end{table}


% ============================================================
% 7. RISK ASSESSMENT
% ============================================================
\section{Risk Assessment}

\begin{table}[h!]
\centering
\small
\renewcommand{\arraystretch}{1.3}
\begin{tabular}{>{\raggedright}p{3.8cm} c >{\raggedright\arraybackslash}p{6.5cm}}
  \toprule
  \textbf{Risk} & \textbf{Likelihood} & \textbf{Mitigation} \\
  \midrule
  SE(3) equivariance broken by spectral coarsening &
  Medium &
  Use center-of-mass aggregation (commutes with rotations/translations);
  add equivariance unit tests at every resolution level. \\

  Non-isotropic posterior computation too expensive for large molecules &
  Medium &
  Lanczos algorithm provides $O(Nk)$ low-rank approximation; fall back to
  block-diagonal approximation if needed. \\

  OMol25 pre-training does not transfer to UV-absorber domain &
  Low &
  OMol25 covers broad chemical diversity including conjugated organics;
  supplement with domain-specific augmentation if transfer is weak. \\

  TD-DFT predictions inaccurate for novel scaffolds &
  Medium &
  Use CAM-B3LYP (validated for charge-transfer states); cross-validate
  against experimental data for known UV absorbers; escalate to
  CC2/ADC(2) for outliers. \\

  Generated molecules not synthesizable &
  Medium &
  Enforce SA score filtering ($< 4.0$) and retrosynthetic feasibility
  during screening; involve synthetic chemist in candidate selection. \\

  Insufficient GPU resources for OMol25-scale pre-training &
  Low &
  Start with 1M-molecule subset; scale up as resources allow.
  Pre-training is parallelizable across multiple nodes. \\

  Photostability of novel absorbers insufficient &
  Medium &
  Include photostability as explicit conditioning signal; screen for
  known photodegradation pathways (e.g., Norrish reactions) computationally
  before synthesis. \\
  \bottomrule
\end{tabular}
\caption{Key project risks and mitigation strategies.}
\label{tab:risks}
\end{table}


% ============================================================
% 8. EXPECTED OUTCOMES AND IMPACT
% ============================================================
\section{Expected Outcomes and Impact}

\subsection{Scientific Contributions}

\begin{enumerate}[leftmargin=*]
  \item \textbf{A unified multi-scale theory for molecular generation.}
    \molssd{} will establish the first rigorous connection between scale-space
    theory~\citep{lindeberg1994scalespace}, molecular
    coarse-graining~\citep{noid2008cg}, and diffusion-based
    generation~\citep{ho2020ddpm}. This theoretical framework --- including the
    information-theoretic justification for spectral coarsening, the non-isotropic
    posterior derivation, and the connection to the renormalization
    group~\citep{wilson1983renormalization} --- will be broadly applicable beyond
    UV absorbers to any molecular generation task.

  \item \textbf{Interpretable multi-scale generation.}
    By producing chemically meaningful intermediate states (chromophore skeleton
    $\to$ fragments $\to$ functional groups $\to$ full atoms), \molssd{} will enable
    a new paradigm of interpretable molecular generation where chemists can understand,
    inspect, and intervene in the design process at multiple levels of abstraction.

  \item \textbf{Computational efficiency gains.}
    We expect multi-resolution generation to achieve 3--5$\times$ wall-clock speedup
    over fixed-resolution baselines on molecules with $>$50 heavy atoms, with larger
    gains for bigger molecules. This will be rigorously benchmarked.

  \item \textbf{Formal bridge between physical and generative coarse-graining.}
    The mathematical equivalence between the SSD forward process on molecular graphs
    and systematic coarse-graining will provide new theoretical tools for both the
    molecular simulation and generative modeling communities.
\end{enumerate}

\subsection{Practical Deliverables}

\begin{enumerate}[leftmargin=*]
  \item \textbf{Novel UV-absorbing molecules.}
    We aim to identify 5--10 novel molecular scaffolds with experimentally validated
    broad-spectrum UV absorption (280--400\,nm), high molar absorptivity
    ($\varepsilon_\text{max} > 10^4$\,L\,mol$^{-1}$\,cm$^{-1}$), and photostability
    exceeding 500 hours of accelerated aging.

  \item \textbf{Computational screening pipeline.}
    A validated, open-source pipeline from conditional molecular generation through
    eSEN energy filtering, synthesizability screening, and TD-DFT spectral
    validation --- reusable for any molecular materials design campaign.

  \item \textbf{UV-absorber dataset.}
    A curated dataset of $\sim$50K UV-absorbing molecules with computed absorption
    spectra, photostability indicators, and physicochemical properties --- a community
    resource for UV-absorber research.

  \item \textbf{Open-source \molssd{} code.}
    A well-documented, modular implementation of \molssd{} with pre-trained weights,
    enabling the community to apply multi-scale diffusion to their own molecular
    design problems.
\end{enumerate}

\subsection{Broader Impact}

\begin{enumerate}[leftmargin=*]
  \item \textbf{Accelerating materials discovery.}
    The \molssd{} framework demonstrates a general paradigm --- principled multi-scale
    generative modeling --- that can accelerate the discovery of functional molecular
    materials across domains: organic semiconductors, photovoltaic materials,
    fluorescent probes, and catalysts.

  \item \textbf{Protecting autonomous systems.}
    Novel UV absorbers with superior performance and photostability will directly
    address the UV degradation challenges facing autonomous vehicle fleets, drones,
    and outdoor robotics, potentially extending component lifetimes by 2--5$\times$.

  \item \textbf{Reducing environmental impact.}
    By designing UV absorbers with reduced migration and leaching, and by
    computationally screening for environmental persistence before synthesis, this
    work contributes to the development of safer, more sustainable protective
    coatings.

  \item \textbf{Advancing multi-scale AI.}
    The theoretical insights from adapting scale-space theory to molecular graphs
    may inspire analogous multi-scale approaches in other structured domains:
    protein design, crystal structure prediction, and polymer engineering.
\end{enumerate}


% ============================================================
% 9. BUDGET CONSIDERATIONS
% ============================================================
\section{Budget Considerations}

\begin{table}[h!]
\centering
\small
\renewcommand{\arraystretch}{1.2}
\begin{tabular}{lrr}
  \toprule
  \textbf{Item} & \textbf{Estimated Cost} & \textbf{Notes} \\
  \midrule
  \multicolumn{3}{l}{\textit{Computational Resources}} \\
  \quad OMol25 pre-training (1M subset)  & \$8,000--12,000  & $\sim$2,000 A100-GPU-hours \\
  \quad QM9/GEOM-Drugs training \& eval  & \$2,000--4,000   & $\sim$500--1,000 GPU-hours \\
  \quad UV-absorber conditional training  & \$4,000--6,000   & $\sim$1,000--1,500 GPU-hours \\
  \quad Generation \& eSEN screening     & \$1,000--2,000   & 10K molecules + inference \\
  \quad TD-DFT (training set, 50K)       & \$15,000--25,000  & $\sim$25K GPU-hours (GPU-accelerated DFT) \\
  \quad TD-DFT (validation, 500)         & \$2,000--4,000   & Higher-quality basis set \\
  \midrule
  \multicolumn{3}{l}{\textit{Wet-Lab Costs}} \\
  \quad Chemical synthesis (10--20 compounds) & \$15,000--30,000 & Reagents, solvents, purification \\
  \quad UV-Vis spectroscopy              & \$2,000--4,000   & Solution and thin-film measurements \\
  \quad Accelerated weathering tests     & \$5,000--10,000  & QUV chamber time, sample preparation \\
  \quad Coating formulation \& testing   & \$5,000--10,000  & Formulation, application, characterization \\
  \midrule
  \multicolumn{3}{l}{\textit{Other}} \\
  \quad Software licenses (Gaussian, etc.)& \$2,000--5,000  & Academic licenses \\
  \quad Publication \& open access fees  & \$5,000--10,000  & Nature OA fee \\
  \midrule
  \textbf{Total}                         & \textbf{\$66,000--122,000} & \\
  \bottomrule
\end{tabular}
\caption{Estimated project budget. GPU costs assume cloud pricing at \$3--5/hour for
A100 instances. Wet-lab costs are order-of-magnitude estimates dependent on molecular
complexity and available infrastructure.}
\label{tab:budget}
\end{table}

\noindent The largest cost drivers are TD-DFT calculations for the training dataset and
wet-lab synthesis. Both can be reduced by (i)~leveraging existing computed spectra from
databases where available, (ii)~using GPU-accelerated DFT codes (e.g., TeraChem, GPU4PySCF)
to reduce per-molecule wall-clock time, and (iii)~prioritizing molecules with short,
well-precedented synthetic routes. The computational budget assumes cloud GPU rental;
costs would be substantially lower if institutional HPC resources are available.


% ============================================================
% 10. REFERENCES
% ============================================================
\newpage
\bibliographystyle{unsrtnat}
\bibliography{references}

\end{document}
